% =====================================================================
%  Partie 3 (~11 min) : ouvrir la boîte -- langage, mémoire, pensée
% =====================================================================
\partie{Ouvrir la boîte}{Ouvrir la boîte : langage, pensée, mémoire}

% ---------------------------------------------------------------------
\begin{frame}{Un cas d'étude : l'accord sujet-verbe}
  \centering
  \begin{tikzpicture}[w/.style={font=\normalsize, text=texte, inner sep=2pt}]
    \node[w] (a) at (0,0) {Les};
    \node[w, text=jaune, font=\normalsize\bfseries, right=0.3cm of a] (b) {clés};
    \node[w, right=0.3cm of b] (c) {que};
    \node[w, right=0.3cm of c] (d) {l'homme};
    \node[w, right=0.3cm of d] (e) {près};
    \node[w, right=0.3cm of e] (f) {de};
    \node[w, right=0.3cm of f] (g) {l'armoire};
    \node[w, right=0.3cm of g] (h) {tient};
    \node[w, text=jaune, font=\normalsize\bfseries, right=0.3cm of h] (i) {sont};
    \node[w, right=0.3cm of i] (j) {\dots};
    \visible<2->{\draw[-{Stealth}, jaune, line width=1.4pt] (b.north) to[out=30, in=150] node[above, font=\scriptsize] {accord à longue distance} (i.north);}
    \visible<3->{\draw[-{Stealth}, rouge, line width=1pt, dashed] (g.south) to[out=-40, in=-140] node[below, font=\scriptsize, text=rouge] {« attracteur » au singulier} (i.south);}
  \end{tikzpicture}

  \vspace{-0.6em}
  \begin{columns}[T,onlytextwidth]
    \column{.48\textwidth}
    \small
    \begin{itemize}
      \item<2-> Le verbe s'accorde avec « clés », \textbf{pas} avec le mot le plus proche.
      \item<3-> L'accord suit la \textbf{structure} de la phrase, pas l'ordre des mots.
    \end{itemize}
    \column{.48\textwidth}
    \small
    \begin{itemize}
      \item<4-> La bonne réponse est connue : même question pour un \textbf{humain} et un \textbf{modèle}.
      \item<5-> Un modèle qui n'a vu que des phrases apprend-il la structure ? \textbf{Comment} ?
    \end{itemize}
  \end{columns}
  \vspace{-0.2em}
  \visible<6->{\cle{Un petit phénomène bien choisi pour tester une grande question.}}
  \raggedright
  \src{Linzen, Dupoux \& Goldberg, TACL 2016}
\end{frame}

% ---------------------------------------------------------------------
\begin{frame}{Dans le modèle : des « neurones du nombre »}
  \begin{columns}[c,onlytextwidth]
    \column{.54\textwidth}
    \centering
    \begin{tikzpicture}[x=0.95cm, y=1.0cm]
      \foreach \w/\i in {Les/0, clés/1, de/2, la/3, porte/4, sont/5}{\node[font=\scriptsize, text=texte] at (\i,-0.35) {\w};}
      \draw[gris] (-0.3,0) -- (5.4,0);
      \draw[orange, line width=1.6pt] plot[smooth] coordinates {(0,0.2) (1,2.0) (2,1.95) (3,1.9) (4,1.92) (5,1.85)};
      \draw[bleu, line width=1.2pt, dashed] plot[smooth] coordinates {(0,0.2) (1,0.15) (2,0.2) (3,0.25) (4,0.2) (5,0.2)};
      \node[font=\scriptsize, text=orange, anchor=west] at (5.5,1.85) {unité « pluriel »};
      \node[font=\scriptsize, text=bleu, anchor=west] at (5.5,0.25) {unité « singulier »};
    \end{tikzpicture}\\[0.2em]
    {\scriptsize\color{gris}schéma d'après Lakretz et al. 2019}
    \column{.44\textwidth}
    \small
    \begin{itemize}
      \item<1-> Dans un réseau LSTM (plus ancien que ChatGPT) de 1\,300 unités, \textbf{deux unités} gardent le nombre du sujet en mémoire.
      \item<2-> Un petit circuit « syntaxique » décide quand les mettre à jour.
      \item<3-> Si on \textbf{supprime} ces unités, le modèle fait des erreurs d'accord.
    \end{itemize}
  \end{columns}
  \vspace{0.4em}
  \visible<4->{\cle{Un mécanisme précis, testable et manipulable, ce qu'on ne peut pas faire dans un cerveau humain.}}
  \src{Lakretz et al., NAACL 2019}
\end{frame}

% ---------------------------------------------------------------------
\begin{frame}{Pensée : le langage n'est pas la pensée}
  \begin{columns}[c,onlytextwidth]
    \column{.44\textwidth}
    \centering
    \cerveau{3.5cm}{%
      \reseaulangage
      \visible<2->{\fill[bleu, opacity=0.75] (0.20,0.78) ellipse (0.07 and 0.06);
      \fill[bleu, opacity=0.75] (0.62,0.80) ellipse (0.09 and 0.06);
      \fill[bleu, opacity=0.6] (0.12,0.60) ellipse (0.04 and 0.05);}}\\[0.2em]
    {\scriptsize \textcolor{rouge}{\textbf{réseau du langage}} \quad \textcolor{bleu}{\textbf{réseau « multi-demande »}}\\\color{gris}(schéma)}
    \column{.54\textwidth}
    \small
    \begin{itemize}
      \item<1-> \textcolor{rouge}{\textbf{Langage}} : comprendre et produire des phrases.
      \item<2-> \textcolor{bleu}{\textbf{Multi-demande}} : calcul, logique, résolution de problèmes.
      \item<3-> Certains patients aphasiques, au langage très atteint, peuvent encore \textbf{calculer et raisonner}.
    \end{itemize}
  \end{columns}
  \vspace{0.5em}
  \visible<4->{\cle{Chez l'humain, langage et pensée s'appuient sur des réseaux en grande partie distincts. Pour les IA, la question reste ouverte.}}
  \src{Fedorenko, Piantadosi \& Gibson, \emph{Nature} 2024}
\end{frame}

% ---------------------------------------------------------------------
\begin{frame}{Mémoire : garder le sujet en tête jusqu'au verbe}
  \begin{columns}[c,onlytextwidth]
    \column{.50\textwidth}
    \centering
    \begin{tikzpicture}[x=0.55cm, y=0.5cm]
      \foreach \i in {0,...,8}{\fill[fondclair] (\i,0) rectangle (\i+0.8,0.8);}
      \fill[jaune] (0,0) rectangle (0.8,0.8);
      \fill[jaune] (8,0) rectangle (8.8,0.8);
      \draw[jaune, line width=1.1pt] (0.4,0.9) to[out=60,in=120] (8.4,0.9);
      \node[font=\scriptsize, text=gris] at (4.4,-0.5) {sujet \dots\ 7 mots \dots\ verbe : coûteux};
      \foreach \i in {0,...,8}{\fill[fondclair] (\i,-2.2) rectangle (\i+0.8,-1.4);}
      \fill[jaune] (3,-2.2) rectangle (3.8,-1.4);
      \fill[jaune] (4,-2.2) rectangle (4.8,-1.4);
      \draw[jaune, line width=1.1pt] (3.4,-1.3) to[out=60,in=120] (4.4,-1.3);
      \node[font=\scriptsize, text=gris] at (4.4,-2.75) {sujet, verbe : facile};
    \end{tikzpicture}
    \column{.48\textwidth}
    \small
    \begin{itemize}
      \item<1-> Notre mémoire de travail est \textbf{limitée} : plus le verbe est loin de son sujet, plus la lecture ralentit.
      \item<2-> Un modèle à qui l'on donne une mémoire \textbf{imparfaite} du contexte reproduit mieux ces difficultés.
    \end{itemize}
  \end{columns}
  \vspace{0.6em}
  \visible<3->{\cle{Si un modèle à la mémoire limitée bute sur les mêmes phrases que nous, c'est que notre difficulté vient de la mémoire.}}
  \src{Gibson, \emph{Cognition} 1998 ; Futrell, Gibson \& Levy, \emph{Cognitive Science} 2020}
\end{frame}
